\documentclass[12pt, a4paper]{article}
\usepackage{graphicx} % Required for inserting images
\graphicspath{{images/}{images-desafios/}}
\usepackage[a4paper, total={7.5in, 10in}]{geometry}
\usepackage{float} % pra conseguir usar o \begin{figure}[H]
\usepackage{indentfirst} % pra fazer tab n 1 paragrafos dos capitulos
% package pra tabela
\usepackage{multirow}
\usepackage{array}

%pacote pra codigo
\usepackage{listings} 
\usepackage{xcolor}
\definecolor{codegreen}{rgb}{0,0.6,0}
\definecolor{codegray}{rgb}{0.5,0.5,0.5}
\definecolor{codepurple}{rgb}{0.58,0,0.82}
\definecolor{backcolour}{rgb}{0.95,0.95,0.92}

\lstdefinestyle{myCodeStyle}{
	backgroundcolor=\color{backcolour},   
	commentstyle=\color{codegreen},
	keywordstyle=\color{cyan},
	numberstyle=\tiny\color{codegray},
	stringstyle=\color{codepurple},
	basicstyle=\ttfamily\footnotesize,
	breakatwhitespace=false,         
	breaklines=true,                 
	captionpos=b,                    
	keepspaces=true,                 
	%numbers=left,                    
	numbersep=5pt,                  
	showspaces=false,                
	showstringspaces=false,
	showtabs=false,                  
	tabsize=2
}
\lstset{style=myCodeStyle}
\renewcommand{\lstlistingname}{Código}%Muda o termo "Listing":  Listing -> Código
\renewcommand{\lstlistlistingname}{Lista de \lstlistingname s}% List of Listings -> Lista de Códigos

% incluir aquivos PDF
\usepackage{pdfpages}


\newcommand{\unit}[1]{\ensuremath{\, \mathrm{#1}}} % comando pra tirar do math mode a unidade de medida

\def\tamanhoGrafico{0.8}

\newcommand{\ohm}{\unit{\Omega} }
\newcommand{\kohm}{\unit{k\Omega}}

\title{
	Laboratório Digital A - PCS3335
	\\  Relatório da Experiência 11 - Projeto da disciplina: Calculadora
}

\author{
	Alunos:
	\\ Rafael Hipólito Rodrigues - N.USP: 14604308
	\\ {{PERSON-2}} - N.USP: {{PERSON-ID-2}}
	\\ \\
	Turma 3 - Bancada A6
	\\ Professor: Bruno de Carvalho Albertini
}

%\date{}

\begin{document}
	
	\maketitle
	\newpage
	
	\tableofcontents
	\newpage
	
	\section{Introdução}
	
	O projeto realizado consiste no desenvolvimento do código VHDL de uma calculadora e a sua aplicação em uma FPGA.
	A calculadora tem funcionamento baseado em Notação Polonesa Inversa e possui uma memória. 
	Foi utilizado um teclado de membrana matricial 4x4 como teclado da calculadora para a sua utilização.
	Os requisito para o projeto da disciplina eram: utilizar uma memória, ter uma Maquina de estados e ter um circuito combinatório não trivial.
	Portando, todos os requisitos foram alcançados com o projeto realizado.
	A figura \ref{fig:fpga} mostra a Calculadora implementada na FPGA.
	
	\begin{figure}[H]
		\centering
		\includegraphics[width=0.6\linewidth]{images/fpga_modo3.jpg}
		\caption{Calculadora implementada na FPGA.}
		\label{fig:fpga}
	\end{figure}
	
	\section{Noções gerais sobre a calculadora}
	\label{sec:nocoes}
	
	Agora será detalhado características gerais sobre a calculadora.
	
	\subsection{Memória}
	
	A memória armazena números de -99.999 até 99.999. A lógica da memória funciona como uma pilha: cada novo número armazenado na memória é inserido no topo
	da pilha, e os números são removidos é retirado do topo da pilha. Essa lógica de funcionamento ajuda a implementação da Notação Polonesa Inversa.
	\par
	Os números são armazenados na memória no formato BCD (\textit{Binary Coded Decimal}). Esse formato foi utilizado pois facilita muito para exibir valores
	decimais em displays de 7 segmentos.
	A memória implementada tem comprimento de 16 palavras, sendo cada palavra de 21 bits. O ultimo bit da palavra armazena no sinal: valor 1 para sinal
	negativo, valor 0 para sinal positivo. Os 20 bits restantes armazenam os 5 digitos decimais em BCD. Por exemplo, o número -09.508 será a palavra:
	
	\begin{table}[H]
		\centering
		\begin{tabular}{c|cccc|cccc|cccc|cccc|cccc}
			Sinal                   & \multicolumn{4}{c|}{Digito 4}                                                & \multicolumn{4}{c|}{Digito 3}                                                & \multicolumn{4}{c|}{Digito 2}                                                & \multicolumn{4}{c|}{Digito 1}                                                & \multicolumn{4}{c}{Digito 0}                                                                      \\ \hline
			\multicolumn{1}{|c|}{1} & \multicolumn{1}{c|}{0} & \multicolumn{1}{c|}{0} & \multicolumn{1}{c|}{0} & 0 & \multicolumn{1}{c|}{1} & \multicolumn{1}{c|}{0} & \multicolumn{1}{c|}{0} & 1 & \multicolumn{1}{c|}{0} & \multicolumn{1}{c|}{1} & \multicolumn{1}{c|}{0} & 1 & \multicolumn{1}{c|}{0} & \multicolumn{1}{c|}{0} & \multicolumn{1}{c|}{0} & 0 & \multicolumn{1}{c|}{1} & \multicolumn{1}{c|}{0} & \multicolumn{1}{c|}{0} & \multicolumn{1}{c|}{0} \\ \hline
		\end{tabular}
		\caption{Valor binário de como o número -09.508 é armazenado.}
	\end{table}
	
	\subsection{Notação Polonesa Inversa}
	
	A lógica de Notação Polonesa Inversa foi escolhida para poder-se realizar expressões matemáticas na calculadora sem o uso de parenteses para indicar
	precedência entre as operações. Por exemplo, a expressão
	\[
	((5+2) \times 6)-(34/(8+7))
	\]
	em Notação Polonesa Inversa é a sequência:
	\[
	5, 2, +, 6, \times, 34, 8, 7, +, /, -
	\]
	
	A lógica pode ser melhor compreendida pelo diagrama da figura \ref{fig:diag_pilha}.
	
	\begin{figure}[H]
		\centering
		\includegraphics[width=1\linewidth]{images/diagrama_pilha.pdf}
		\caption{Diagrama que mostra a evolução da pilha no exemplo.}
		\label{fig:diag_pilha}
	\end{figure}
	
	É importante notar que as operações de subtração e divisão são realizados seguindo a lógica:
	
	\begin{itemize}
		\item Subtração: [Valor antes do topo da pilha] - [Valor no topo da pilha]
		
		\item Divisão: [Valor antes do topo da pilha] / [Valor no topo da pilha]
	\end{itemize}
	
	Observe que a operação de divisão resulta apenas na parte inteira do resultado.
	
	\subsection{Modos de funcionamento}
	
	A calculadora é dividida em 3 modos de funcionamento para facilitar a usabilidade.
	
	\subsubsection*{Modo 1: armazenamento de números na memória}
	No 1° modo o usuário apenas armazena números na memória, adicionando valores ao topo da pilha.
	
	\subsubsection*{Modo 2: realização de operações aritméticas com os valores da pilha}
	No 2° modo o usuário apenas realiza operações aritméticas com os valores que já estão na armazenados na memoria. Ao selecionar 
	uma operação aritméticas, a calculadora pegar os 2 números mais ao topo da pilha, realizaria a operação entre eles, remove esses 2 números da 
	pilha, e armazenaria no topo da pilha o resultado da operação.
	
	\subsubsection*{Modo 3: exibição dos valores da pilha}
	No 3° modo o usuário poderá ler os valores que ja estão armazenados na memoria.
	
	\subsection{Outras características}
	
	\begin{itemize}
		\item Somente é possível inserir e realizar operações aritméticas com números reais inteiros com sinal.
		
		\item A operação de divisão resulta no valor inteiro da divisão.
		
		\item A calculadora não trata de casos de overflow nas operações aritméticas. Caso o resultado de alguma operação ultrapasse o valor máximo (-99.999 até 99.999),
		um valor incorreto será armazenado na memória.
		
		\item O uso da calculadora é realizado através de um teclado de membrana matricial 4x4 e das chaves da FPGA, como mostrado na figura \ref{fig:fpga}.
	\end{itemize}
	
	\section{Manual}
	
	Nessa seção será detalhado como utilizar a calculadora tendo ela implementada numa FPGA com o teclado de membrana matricial 4x4.
	
	\subsection{Menu inicial}
	
	Ao iniciar a calculadora, ela estará no menu inicial, onde exibirá "SELECT" nos displays de 7 segmentos, assim como mostrado a figura \ref{fig:fpga_menu}.
	A parti desse menu é possível ir para os modos 1, 2 e 3. Para ir ao modo desejado, basta aperta o número do modo no teclado.
	\par
	Estando em um dos modos de funcionamento, para voltar ao menu inicial coloque o chave 0 da FPGA (SW0) no valor Lógico 1. Ao sair de um modo pela chave 0, o modo 
	não é reiniciado. Ou seja, ao sair de um modo e voltar para ele, ele continuará na mesma situação ao ter saído dele.
	\par
	A chave 9 da FPGA (SW1), ao ser ativada reseta a calculadora, reiniciando todos so modos (colocando todos seus contadores e registradores no valor inicial) e 
	vai para o menu inicial. Contudo, a memoria não é reiniciada ao ativar a chave 9.
	
	\begin{figure}[H]
		\centering
		\includegraphics[width=0.6\linewidth]{images/fpga_menu.jpg}
		\caption{Calculadora no menu inicial.}
		\label{fig:fpga_menu}
	\end{figure}
	
	\subsection{Modo 1}
	
	O modo 1 permite o usuário digitar valores e armazenar eles na memória.
	Ao iniciar o modo 1, o usuário pode aperta digitar o número desejado utilizando o teclado de membrana. O valor digitado vai sendo exibido nos displays de 7
	segmentos. O display mais a esquerda mostra o sinal do número. Os 5 displays mais a direita mostram o valor do número.
	\par
	Os dígitos apertados no teclado serão alocados do menos significativo para o mais significativo no número. 
	O display que mostra "\_", indica que o próximo digito a ser
	teclado entrará naquela posição e atualmente o valor do digito nessa posição é 0. 
	A figura \ref{fig:fpga_modo1} mostra o modo 1 em funcionamento. 
	
	\begin{figure}[H]
		\centering
		\includegraphics[width=0.6\linewidth]{images/fpga_modo1.jpg}
		\caption{Calculadora no modo 1 de funcionamento como o valor -00502 digitado.}
		\label{fig:fpga_modo1}
	\end{figure}
	
	O botão "C" no teclado apaga o ultimo digito do número.
	O botão "F1" no teclado inverte o sinal do número.
	O botão "E" no teclado armazena o valor digitado na memória, colocando-o no topo da pilha, e zera o valor mostrado nos displays.
	\par
	Não é possível digitar um número maior que 5 dígitos. Não é possível utilizar o botão apagar quando há nenhum digito no número.
	
	\subsection{Modo 2}
	
	O modo 2 permite que o usuário realize operações aritméticas com os valores armazenados na memória.
	Nesse modo, os displays de 7 seg. sempre exibiram o número no topo da pilha.
	As operações aritméticas são realizadas com os 2 números mais ao topo da pilha, como explicado na seção \ref{sec:nocoes}.
	Para realizar as operações usa-se os botões:
	\begin{itemize}
		\item F1: adição
		\item F2: subtração
		\item F3: divisão
		\item F4: multiplicação
	\end{itemize}
	As operações aritméticas somente serão realizadas quando houver pelo menos 2 números armazenados na memória.
	
	\subsection{Modo 3}
	
	O modo 3 permite que o usuário veja os valores já armazenados na memória. Ao iniciar o modo 3, ele exibe nos displays de 7 seg. o número armazenado na 
	posição 0 da memória. 
	Ao aperta os botões:
	\begin{itemize}
		\item 8: a posição do número que será mostrado nos display é incrementado em 1. Ou seja, se o número da posição 3 era exibido, ao pressionar "8", será exibido o
		número da posição 4.
		\item 2: a posição do número que será mostrado nos display é diminuído em 1. Ou seja, se o número da posição 3 era exibido, ao pressionar "2", será exibido o
		número da posição 2.
	\end{itemize}
	Lembrando que a memória consegue armazenar 16 números, que vão da posição 0 até a posição 15.
	A figura \ref{fig:fpga_modo3} mostra a calculadora no modo 3 exibindo o número armazenado na posição 1 da memória.
	
	\begin{figure}[H]
		\centering
		\includegraphics[width=0.6\linewidth]{images/fpga_modo3.jpg}
		\caption{Calculadora no modo 3 de funcionamento.}
		\label{fig:fpga_modo3}
	\end{figure}
	
	Importante destacar que esse modo funciona independente da pilha. Esse modo mostra todas as 16 posições da memória, independente de quanto valores há na pilha.
	Assim, mostrando valores que não estão na pilha. Contudo, o valor inicial para todas as posições da memória é zero. 
	E também, ao remover um elemento da pilha, a posição da memória que esse elemento ocupava vai ficar com valor zero.
	Portando, posições da memória que não estão sendo utilizadas na pilha sempre exibirão valor zero.
	\par
	Os leds nesse modo mostram:
	\begin{itemize}
		\item Os 4 primeiros leds da FPGA (leds 0, 1, 2 e 3) exibem o valor binário da posição do número que está sendo atualmente exibido nos displays.
		\item Os 4 próximos leds (leds 4, 5, 6 e 7) exibem o valor binário de quantos números há armazenados na pilha.
	\end{itemize}
	Por exemplo, na figura \ref{fig:fpga_modo3}, os leds mostram que há 3 números na pilha e o valor atualmente mostrado no display é a posição 1 da memória.
	
	
	\section{Funcionamento do circuito}
	
	Nessa seção será explicado de forma simplificada o funcionamento e o circuito de cada entidade implementada no código VHDL da calculadora.
	O código VHDL foi compilado e implementado na FPGA utilizando o software Quartus.
	A calculadora possui as entidades:
	
	\begin{itemize}
		\item \texttt{ascii2seg}
		\item \texttt{baudRateGenerator}
		\item \texttt{bcd\_asciiConverter}
		\item \texttt{BCD\_to\_bin\_UNSIGNED}
		\item \texttt{BCD\_to\_bin}
		\item \texttt{bin\_to\_BCD\_UNSIGNED}
		\item \texttt{bin\_to\_BCD}
		\item \texttt{calculator} (top-level)
		\item \texttt{counter}
		\item \texttt{keyboard4x4}
		\item \texttt{memory}
		\item \texttt{mode1}
		\item \texttt{mode2}
		\item \texttt{mode3}
		\item \texttt{shiftregister}
		\item \texttt{shiftregisterKeyboard}
	\end{itemize}
	
	As entidades \texttt{ascii2seg}, \texttt{baudRateGenerator}, \texttt{counter} e \texttt{shiftregister} são as mesmas utilizadas nas 
	experiencias anteriores que implementavam o modulo UART.
	\par
	A entidade \texttt{shiftregisterKeyboard} é praticamente idêntica ao \texttt{shiftregister}, a única diferença é que esse registrador
	inicia o com valor "011111..." ao invés de "00000...". Ele é utilizado da entidade texttt{keyboard4x4}.
	
	\subsection{\texttt{calculator}}
	
	A entidade \texttt{calculator} é a top-level. A lógica dela controla qual entidade esta sendo usado no momento.
	
	\begin{figure}[H]
		\centering
		\includegraphics[width=0.7\linewidth]{images/rtl_calculator_zoomin.pdf}
		\caption{RTL da entidade.}
		\label{fig:rtl_calculator_zoomin}
	\end{figure}
	
	\begin{figure}[H]
		\centering
		\includegraphics[width=1\linewidth]{images/rtl_calculadora.png}
		\caption{RTL da entidade.}
		\label{fig:rtl_calculadora}
	\end{figure}
	
	\begin{figure}[H]
		\centering
		\includegraphics[width=0.6\linewidth]{images/fsm_calculadora.png}
		\caption{FSM da entidade.}
		\label{fig:fsm_calculadora}
	\end{figure}
	
	\subsection{\texttt{keyboard4x4}}
	
	A entidade \texttt{keyboard4x4} controla o teclado de membrana matricial 4x4. Ela detecta quando uma tecla é apertada e manda 
	o ASCII dessa teclada para as outras entidades.
	
	\begin{figure}[H]
		\centering
		\includegraphics[width=0.4\linewidth]{images/rtl_keyboard.png}
		\caption{RTL da entidade.}
		\label{fig:rtl_keyboard}
	\end{figure}
	
	\begin{figure}[H]
		\centering
		\includegraphics[width=1\linewidth]{images/rtl_keyboard_zoom.png}
		\caption{RTL da entidade.}
		\label{fig:rtl_keyboard_zoom}
	\end{figure}
	
	\begin{figure}[H]
		\centering
		\includegraphics[width=0.7\linewidth]{images/fsm_keyboard.png}
		\caption{FSM da entidade.}
		\label{fig:fsm_keyboard}
	\end{figure}
	
	\subsection{\texttt{memory}}
	
	A entidade \texttt{memory} instancia a memoria utilizada pela calculadora. Na figura \ref{fig:rtl_memory} é possível observar 
	um bloxo roxo chamado "memoria\_ram". Esse bloco foi a memória que o Quartus instanciou.
	
	\begin{figure}[H]
		\centering
		\includegraphics[width=0.6\linewidth]{images/rtl_memory.png}
		\caption{RTL da entidade.}
		\label{fig:rtl_memory}
	\end{figure}
	
	\subsection{\texttt{mode1}}
	
	A entidade \texttt{mode1} implementa o modo 1 de funcionamento.
	
	\begin{figure}[H]
		\centering
		\includegraphics[width=0.6\linewidth]{images/rtl_mode1.png}
		\caption{RTL da entidade.}
		\label{fig:rtl_mode1}
	\end{figure}
	
	\begin{figure}[H]
		\centering
		\includegraphics[width=1\linewidth]{images/rtl_mode1_zoom.png}
		\caption{RTL da entidade.}
		\label{fig:rtl_mode1_zoom}
	\end{figure}
	
	Seguindo a FSM da entidade na figura \ref{fig:fsm_mode1} é possível realizar 3 ações: Adicionar um digito ao registrador, 
	remover um digito do registrador e armazenar o valor do registrado na memória.
	
	\begin{figure}[H]
		\centering
		\includegraphics[width=1\linewidth]{images/fsm_mode1.png}
		\caption{FSM da entidade.}
		\label{fig:fsm_mode1}
	\end{figure}
	
	
	
	\subsection{\texttt{mode2}}
	
	A entidade \texttt{mode2} implementa o modo 2 de funcionamento.
	
	\begin{figure}[H]
		\centering
		\includegraphics[width=0.5\linewidth]{images/rtl_mode2.png}
		\caption{RTL da entidade.}
		\label{fig:rtl_mode2}
	\end{figure}
	
	\begin{figure}[H]
		\centering
		\includegraphics[width=1\linewidth]{images/rtl_mode2_zoom.png}
		\caption{RTL da entidade.}
		\label{fig:rtl_mode2_zoom}
	\end{figure}
	
	Seguindo a FSM da entidade na figura \ref{fig:fsm_mode2} é possível perceber a sequencia de passos:
	armazenar temporariamente a operação selecionada, armazenar temporariamente os 2 operandos, transformar o resultado da operação 
	de binário para BCD e atualizar a memória com o valor do resultado.
	\par
	É importante destacar que as operações aritméticas em binário, e a conversão de BCD para binario utiliza circuitos combinatórios, logo 
	não é necessário um estado para esses passos.
	
	
	\begin{figure}[H]
		\centering
		\includegraphics[width=1\linewidth]{images/fsm_mode2.png}
		\caption{FSM da entidade.}
		\label{fig:fsm_mode2}
	\end{figure}
	
	
	\subsection{\texttt{mode3}}
	
	A entidade \texttt{mode3} implementa o modo 3 de funcionamento.
	
	\begin{figure}[H]
		\centering
		\includegraphics[width=0.5\linewidth]{images/rtl_mode3.png}
		\caption{RTL da entidade.}
		\label{fig:rtl_mode3}
	\end{figure}
	
	Na RTL da entidade na figura \ref{fig:rtl_mode3_zoom} é possível que há um contador que armazena o valor da posição que está 
	sendo exibida atualmente nos display. 
	
	\begin{figure}[H]
		\centering
		\includegraphics[width=1\linewidth]{images/rtl_mode3_zoom.png}
		\caption{RTL da entidade.}
		\label{fig:rtl_mode3_zoom}
	\end{figure}
	
	Seguindo a FSM da entidade na figura \ref{fig:fsm_mode3}: de acordo com a tecla apertada, o valor do contador mencionado anteriormente 
	aumenta ou diminui em uma unidade.
	
	\begin{figure}[H]
		\centering
		\includegraphics[width=0.9\linewidth]{images/fsm_mode3.png}
		\caption{FSM da entidade.}
		\label{fig:fsm_mode3}
	\end{figure}
	
	
	
	\subsection{\texttt{BCD\_to\_bin\_UNSIGNED}}
	
	A entidade \texttt{BCD\_to\_bin\_UNSIGNED} é um circuito combinatório que transforma BCD de 21 bits em binário de 18 bits SEM SINAL.
	
	\begin{figure}[H]
		\centering
		\includegraphics[width=0.6\linewidth]{images/rtl_bcdToBin_un.png}
		\caption{RTL da entidade.}
		\label{fig:rtl_bcdToBin_un}
	\end{figure}
	
	
	
	\subsection{\texttt{BCD\_to\_bin}}
	
	A entidade \texttt{BCD\_to\_bin} é o circuito combinatório que acrescenta a lógica de sinal (complemento de 2) à 
	saída do \texttt{BCD\_to\_bin\_UNSIGNED}.
	
	\begin{figure}[H]
		\centering
		\includegraphics[width=0.6\linewidth]{images/rtl_bcdToBin.png}
		\caption{RTL da entidade.}
		\label{fig:rtl_bcdToBin}
	\end{figure}
	
	
	
	\subsection{\texttt{bin\_to\_BCD\_UNSIGNED}}
	
	A entidade \texttt{bin\_to\_BCD\_UNSIGNED} implementa o algoritomo Double dabble para converter de binario para BCD.
	Os valores binarios na entrada da entidade são binarios SEM SINAL.
	
	\begin{figure}[H]
		\centering
		\includegraphics[width=0.6\linewidth]{images/rtl_binToBcd_un_zoomin.png}
		\caption{RTL da entidade.}
		\label{fig:rtl_binToBcd_un_zoomin}
	\end{figure}
	
	\begin{figure}[H]
		\centering
		\includegraphics[width=0.6\linewidth]{images/rtl_binToBcd_un.png}
		\caption{RTL da entidade.}
		\label{fig:rtl_binToBcd_un}
	\end{figure}
	
	\begin{figure}[H]
		\centering
		\includegraphics[width=0.8\linewidth]{images/fsm_binToBCD.png}
		\caption{FSM da entidade.}
		\label{fig:fsm_binToBCD}
	\end{figure}
	
	
	
	\subsection{\texttt{bin\_to\_BCD}}
	
	A entidade \texttt{bin\_to\_BCD} instancia a entidade \texttt{bin\_to\_BCD\_UNSIGNED}. Ela realiza a lógica de transformar 
	um binário COM SINAL em um SEM SINAL, passar esse novo binaria para a entidade \texttt{bin\_to\_BCD\_UNSIGNED}, pegar a 
	saída BCD SEM SINAL de \texttt{bin\_to\_BCD\_UNSIGNED} e aplicar o sinal a esse BCD conforme o binário COM SINAL original.
	
	\begin{figure}[H]
		\centering
		\includegraphics[width=0.8\linewidth]{images/rtl_binToBcd.png}
		\caption{RTL da entidade.}
		\label{fig:rtl_binToBcd}
	\end{figure}

	
	\section{Implementação e Testes}
	
	A implementação da calculadora na FPGA realizada pelo grupo utilizou a seguinte pinagem:
	
	\begin{table}[H]
		\begin{center}
			\begin{tabular}{|l|l|l|}
				\hline
				c{[}3{]}            & PIN\_P16  & Output \\ \hline
				c{[}2{]}            & PIN\_L22  & Output \\ \hline
				c{[}1{]}            & PIN\_M22  & Output \\ \hline
				c{[}0{]}            & PIN\_N20  & Output \\ \hline
				clock50M            & PIN\_M9   & Input  \\ \hline
				display7seg1{[}6{]} & PIN\_AA22 & Output \\ \hline
				... pinos dos display de 7 seg. da FPGA ... & ...       & Output \\ \hline
				display7seg6{[}0{]} & PIN\_N9   & Output \\ \hline
				goToMenu            & PIN\_U13  & Input  \\ \hline
				l{[}3{]}            & PIN\_L19  & Input  \\ \hline
				l{[}2{]}            & PIN\_L18  & Input  \\ \hline
				l{[}1{]}            & PIN\_N21  & Input  \\ \hline
				l{[}0{]}            & PIN\_R21  & Input  \\ \hline
				leds{[}9{]}         & PIN\_L2   & Output \\ \hline
				... pinos dos leds da FPGA ...  & ...       & Output \\ \hline
				leds{[}0{]}         & PIN\_AA2  & Output \\ \hline
				pinDebug            & PIN\_N16  & Output \\ \hline
				reset               & PIN\_AB12 & Input  \\ \hline
			\end{tabular}
		\end{center}
	\end{table}
	
	\subsection{Testes}
	
	O grupo testou os 3 modos de funcionamento e todos apresentaram funcionamento correto.
	
	Observações: o grupo percebeu certa instabilidade no teclado de membrana matricial utilizado.
	Em algumas situações ele não respondia corretamente, gerando um funcionamento incorreto na calculadora.
	
	\section{Anexos}
	
	Nas paginas a seguir é apresentado o RTL de algumas entidades em imagens vetoriais para a melhor visualização.
	
	\includepdf[pages=-]{images/rtl_calculadora_pdf.pdf}
	
	\includepdf[]{images/rtl_keyboard_pdf.pdf}
	
	\includepdf[]{images/rtl_mode1_pdf.pdf}
	
	\includepdf[]{images/rtl_mode2_pdf.pdf}

	
	
\end{document}
